\chapter{Chapter07 — Quadratic PCPP}
%%% generated by the blueprint pipeline from FABL.Chapter07.QuadraticPCPP
\section{Overview}

Declarations of the Lean module \texttt{FABL.Chapter07.QuadraticPCPP}.

\section{Declarations}

\begin{definition}[Flattening a binary matrix]
\lean{FABL.flattenBinaryMatrix}
\label{FABL.flattenBinaryMatrix}
\leanok
\uses{FABL.F₂Cube, FABL.𝔽₂}
For a natural number $n$ and an $n\times n$ matrix $A$ over the field with two elements,
flatten $A$ into a vector of length $n^2$ over that field by listing its entries in the
order induced by the standard equivalence between pairs of indices and indices in
$	ext{Fin}(n^2)$.
\end{definition}

\begin{definition}[Reconstructing a binary matrix]
\lean{FABL.unflattenBinaryMatrix}
\label{FABL.unflattenBinaryMatrix}
\leanok
\uses{FABL.F₂Cube, FABL.𝔽₂}
For a natural number $n$ and a vector $x$ of length $n^2$ over the field with two
elements, define the $n\times n$ matrix whose $(i,j)$ entry is the coordinate of $x$
indexed by the standard encoding of the pair $(i,j)$.
\end{definition}

\begin{theorem}[Flattening After Unflattening]
\lean{FABL.flattenBinaryMatrix_unflattenBinaryMatrix}
\label{FABL.flattenBinaryMatrix_unflattenBinaryMatrix}
\leanok
\uses{FABL.F₂Cube, FABL.flattenBinaryMatrix, FABL.unflattenBinaryMatrix}
\difficulty{3}
For a natural number $n$, flattening the matrix obtained by unflattening any vector of
length $n^2$ over the field with two elements returns the original vector.
\end{theorem}

\begin{theorem}[Unflattening a flattened matrix]
\lean{FABL.unflattenBinaryMatrix_flattenBinaryMatrix}
\label{FABL.unflattenBinaryMatrix_flattenBinaryMatrix}
\leanok
\uses{FABL.flattenBinaryMatrix, FABL.unflattenBinaryMatrix, FABL.𝔽₂}
\difficulty{3}
For every natural number $n$ and every $n 	imes n$ matrix $A$ over $\F_2$, unflattening
its associated vector recovers the matrix:
$\operatorname{unflatten}(\operatorname{flatten}(A))=A$.
\end{theorem}

\begin{definition}[Binary matrices and binary vectors]
\lean{FABL.binaryMatrixCubeEquiv}
\label{FABL.binaryMatrixCubeEquiv}
\leanok
\uses{FABL.F₂Cube, FABL.flattenBinaryMatrix, FABL.flattenBinaryMatrix_unflattenBinaryMatrix, FABL.unflattenBinaryMatrix, FABL.unflattenBinaryMatrix_flattenBinaryMatrix, FABL.𝔽₂}
For each natural number $n$, the equivalence between $n\times n$ matrices over the field
with two elements and vectors of length $n^2$ over that field sends a matrix to its
flattened vector and sends a vector to the matrix reconstructed from its coordinates.
\end{definition}

\begin{theorem}[Flattening the zero matrix]
\lean{FABL.flattenBinaryMatrix_zero}
\label{FABL.flattenBinaryMatrix_zero}
\leanok
\uses{FABL.flattenBinaryMatrix, FABL.𝔽₂}
\difficulty{1}
For every natural number $n$, flattening the zero $n\times n$ matrix over the field with
two elements gives the zero vector in the $n^2$-dimensional Boolean cube.
\end{theorem}

\begin{theorem}[Flattening Preserves Matrix Addition]
\lean{FABL.flattenBinaryMatrix_add}
\label{FABL.flattenBinaryMatrix_add}
\leanok
\uses{FABL.flattenBinaryMatrix, FABL.𝔽₂}
\difficulty{1}
For a natural number $n$ and two $n\times n$ matrices over the field with two elements,
flattening their sum into a vector of length $n^2$ gives the sum of their flattened
vectors.
\end{theorem}

\begin{theorem}[Unflattening the zero vector]
\lean{FABL.unflattenBinaryMatrix_zero}
\label{FABL.unflattenBinaryMatrix_zero}
\leanok
\uses{FABL.F₂Cube, FABL.unflattenBinaryMatrix}
\difficulty{1}
For every natural number $n$, unflattening the zero vector of length $n^2$ over $\F_2$
gives the zero $n 	imes n$ matrix.
\end{theorem}

\begin{theorem}[Additivity of matrix unflattening]
\lean{FABL.unflattenBinaryMatrix_add}
\label{FABL.unflattenBinaryMatrix_add}
\leanok
\uses{FABL.F₂Cube, FABL.unflattenBinaryMatrix}
\difficulty{1}
For every natural number $n$ and all vectors $x,y \in (\F_2)^{n^2}$, unflattening their
coordinatewise sum gives the sum of their associated $n 	imes n$ matrices over $\F_2$:
$\operatorname{unflatten}(x+y)=\operatorname{unflatten}(x)+\operatorname{unflatten}(y)$.
\end{theorem}

\begin{definition}[Binary matrix pairing]
\lean{FABL.binaryMatrixPairing}
\label{FABL.binaryMatrixPairing}
\leanok
\uses{FABL.𝔽₂}
For two $n\times n$ matrices $A$ and $B$ over the field with two elements, their pairing
is the sum, over all row and column indices, of the products of their corresponding
entries.
\end{definition}

\begin{definition}[Binary outer product]
\lean{FABL.binaryOuterProduct}
\label{FABL.binaryOuterProduct}
\leanok
\uses{FABL.F₂Cube, FABL.𝔽₂}
For vectors $x,y$ of length $n$ over the field with two elements, their outer product is
the $n\times n$ matrix whose $(i,j)$ entry is $x_i y_j$.
\end{definition}

\begin{definition}[Binary matrix character]
\lean{FABL.binaryMatrixCharacter}
\statementsource{Analysis-of-Boolean-Functions-by-Ryan-ODonnell}{
  pdf-ffbe9c7b3306-p023-b001,
  pdf-ffbe9c7b3306-p071-b009
}
\label{FABL.binaryMatrixCharacter}
\leanok
\uses{FABL.binaryMatrixPairing, FABL.𝔽₂}
For a binary matrix $A$, the associated binary matrix character is the function sending
each binary matrix $B$ to the binary matrix pairing of $A$ and $B$, viewed as an element
of $\mathbb{F}_2$.
\end{definition}

\begin{theorem}[Dot product of flattened binary matrices]
\lean{FABL.f₂DotProduct_flattenBinaryMatrix}
\leanok
\uses{FABL.binaryMatrixPairing, FABL.flattenBinaryMatrix, FABL.f₂DotProduct, FABL.𝔽₂}
\difficulty{4}
Let $n$ be a nonnegative integer and let $A,B$ be $n\times n$ matrices over $\bbr_2$.
Flatten each matrix into a vector using the coordinate ordering of pairs of indices. The
dot product of the flattened matrices equals $\sum_{i=0}^{n-1}\sum_{j=0}^{n-1}
A_{ij}B_{ij}$ in $\bbr_2$.
\end{theorem}

\begin{theorem}[Pairing with an Outer Product]
\lean{FABL.binaryMatrixPairing_rankOne}
\label{FABL.binaryMatrixPairing_rankOne}
\leanok
\uses{FABL.F₂Cube, FABL.binaryMatrixPairing, FABL.binaryOuterProduct, FABL.matrixRankOnePairing, FABL.𝔽₂}
\difficulty{1}
Let $n$ be a natural number, let $A$ be an $n\times n$ matrix over $\mathbb{F}_2$, and
let $x,y$ be vectors of length $n$ over $\mathbb{F}_2$. The pairing of $A$ with the
outer product whose $(i,j)$ entry is $x_i y_j$ equals $\sum_{i,j} A_{ij}x_i y_j$ in
$\mathbb{F}_2$.
\end{theorem}

\begin{theorem}[Additivity of the matrix character]
\lean{FABL.binaryMatrixCharacter_add}
\label{FABL.binaryMatrixCharacter_add}
\leanok
\uses{FABL.binaryMatrixCharacter, FABL.𝔽₂}
\difficulty{3}
Let $n$ be a natural number and let $A,B,C$ be $n\times n$ matrices over $\mathbb F_2$.
The sum of the entrywise products of $A$ with $B+C$ equals the sum of the entrywise
products of $A$ with $B$, plus the sum of the entrywise products of $A$ with $C$.
\end{theorem}

\begin{theorem}[Matrix Character Under Unflattening]
\lean{FABL.binaryMatrixCharacter_unflattenBinaryMatrix}
\label{FABL.binaryMatrixCharacter_unflattenBinaryMatrix}
\leanok
\uses{FABL.F₂Cube, FABL.binaryMatrixCharacter, FABL.flattenBinaryMatrix, FABL.f₂DotProduct, FABL.f₂DotProduct_flattenBinaryMatrix, FABL.unflattenBinaryMatrix, FABL.𝔽₂}
\difficulty{3}
Let $n$ be a natural number, let $A$ be an $n\times n$ matrix over $\mathbb{F}_2$, and
let $x$ be a vector of length $n^2$ over $\mathbb{F}_2$. The matrix character of $A$
evaluated at the matrix obtained by unflattening $x$ equals the $\mathbb{F}_2$ dot
product of the flattening of $A$ with $x$.
\end{theorem}

\begin{definition}[Quadratic character address space]
\lean{FABL.QuadraticCharacterAddress}
\statementsource{Analysis-of-Boolean-Functions-by-Ryan-ODonnell}{
  pdf-ffbe9c7b3306-p197-b012,
  pdf-ffbe9c7b3306-p198-b003
}
\label{FABL.QuadraticCharacterAddress}
\leanok
\uses{FABL.F₂Cube, FABL.𝔽₂}
For $n\ge 0$, a quadratic character address is either a vector in $\mathbb{F}_2^n$ or an
$n\times n$ matrix over $\mathbb{F}_2$.
\end{definition}

\begin{definition}[Quadratic character oracle]
\lean{FABL.quadraticCharacterOracle}
\statementsource{Analysis-of-Boolean-Functions-by-Ryan-ODonnell}{
  pdf-ffbe9c7b3306-p197-b012,
  pdf-ffbe9c7b3306-p198-b003
}
\label{FABL.quadraticCharacterOracle}
\leanok
\uses{FABL.F₂Cube, FABL.QuadraticCharacterAddress, FABL.𝔽₂}
Let $\ell\colon\mathbb{F}_2^n\to\mathbb{F}_2$ and $q\colon\mathbb{F}_2^{n\times
n}\to\mathbb{F}_2$. The oracle on the disjoint union of vectors and binary matrices
returns $\ell(x)$ on a vector address $x$ and $q(M)$ on a matrix address $M$.
\end{definition}

\begin{definition}[Quadratic pair linearity tester]
\lean{FABL.quadraticPairLinearBLRTester}
\statementsource{Analysis-of-Boolean-Functions-by-Ryan-ODonnell}{
  pdf-ffbe9c7b3306-p032-b002,
  pdf-ffbe9c7b3306-p197-b010,
  pdf-ffbe9c7b3306-p197-b012,
  pdf-ffbe9c7b3306-p198-b003
}
\label{FABL.quadraticPairLinearBLRTester}
\leanok
\uses{FABL.F₂Cube, FABL.NonadaptiveFunctionTester, FABL.QuadraticCharacterAddress, FABL.uniformPMF, FABL.𝔽₂}
For $n\ge 0$, choose independent uniform vectors $x,y\in\mathbb{F}_2^n$. The
four-position nonadaptive tester queries an oracle at the vector addresses $x$, $y$, and
$x+y$, leaves its fourth position unanswered, and accepts exactly when the three
returned values $a,b,c$ satisfy $a+b=c$ in $\mathbb{F}_2$; if any of the first three
answers is absent, it rejects.
\end{definition}

\begin{definition}[Quadratic pair matrix tester]
\lean{FABL.quadraticPairMatrixBLRTester}
\statementsource{Analysis-of-Boolean-Functions-by-Ryan-ODonnell}{
  pdf-ffbe9c7b3306-p197-b012,
  pdf-ffbe9c7b3306-p198-b003
}
\label{FABL.quadraticPairMatrixBLRTester}
\leanok
\uses{FABL.F₂Cube, FABL.NonadaptiveFunctionTester, FABL.QuadraticCharacterAddress, FABL.unflattenBinaryMatrix, FABL.uniformPMF, FABL.𝔽₂}
For each $n\ge 0$, consider a tester on addresses that are either vectors in $\mathbb
F_2^n$ or $n\times n$ matrices over $\mathbb F_2$. It samples two independent uniform
binary vectors of length $n^2$, queries the matrices obtained by reshaping each vector
and their sum, and makes no fourth query. It accepts exactly when the first two returned
values sum to the third in $\mathbb F_2$; any missing answer causes rejection.
\end{definition}

\begin{definition}[Quadratic pair rank-one tester]
\lean{FABL.quadraticPairRankOneTester}
\statementsource{Analysis-of-Boolean-Functions-by-Ryan-ODonnell}{
  pdf-ffbe9c7b3306-p198-b001,
  pdf-ffbe9c7b3306-p198-b003
}
\label{FABL.quadraticPairRankOneTester}
\leanok
\uses{FABL.F₂Cube, FABL.NonadaptiveFunctionTester, FABL.QuadraticCharacterAddress, FABL.binaryOuterProduct, FABL.unflattenBinaryMatrix, FABL.uniformPMF, FABL.𝔽₂}
For each $n\ge 0$, consider a tester on addresses that are either vectors in $\mathbb
F_2^n$ or $n\times n$ matrices over $\mathbb F_2$. It samples independent uniform
vectors $x,y\in\mathbb F_2^n$ and $z\in\mathbb F_2^{n^2}$, then queries $x$, $y$, the
matrix obtained by reshaping $z$, and that matrix plus the outer product $xy^{\mathsf
T}$. It accepts exactly when the product of the first two returned values equals the sum
of the last two in $\mathbb F_2$; any missing answer causes rejection.
\end{definition}

\begin{definition}[Quadratic character pair tester]
\lean{FABL.quadraticCharacterPairTester}
\statementsource{Analysis-of-Boolean-Functions-by-Ryan-ODonnell}{
  pdf-ffbe9c7b3306-p198-b001,
  pdf-ffbe9c7b3306-p198-b003
}
\label{FABL.quadraticCharacterPairTester}
\leanok
\uses{FABL.NonadaptiveFunctionTester, FABL.QuadraticCharacterAddress, FABL.quadraticPairLinearBLRTester, FABL.quadraticPairMatrixBLRTester, FABL.quadraticPairRankOneTester, FABL.uniformMixtureFunctionTester, FABL.𝔽₂}
For each $n\ge 0$, consider the tester obtained by choosing uniformly among three
testers: the linear BLR tester, the matrix BLR tester, and the rank-one tester. The
chosen tester is run on addresses that are either vectors in $\mathbb F_2^n$ or $n\times
n$ matrices over $\mathbb F_2$, with answers in $\mathbb F_2$ and four query positions.
\end{definition}

\begin{theorem}[Quadratic-pair tester acceptance probability]
\lean{FABL.quadraticPairLinearBLRTester_acceptanceProbability}
\label{FABL.quadraticPairLinearBLRTester_acceptanceProbability}
\leanok
\uses{FABL.F₂Cube, FABL.NonadaptiveFunctionTester.acceptanceProbability, FABL.blrFunctionTester, FABL.pmfExpectation, FABL.quadraticCharacterOracle, FABL.quadraticPairLinearBLRTester, FABL.uniformPMF, FABL.𝔽₂}
\difficulty{3}
For every natural number $n$, let $\ell:(\mathbb{F}_2^n)\to\mathbb{F}_2$ and let $q$
assign an element of $\mathbb{F}_2$ to each $n\times n$ matrix over $\mathbb{F}_2$. The
probability that the tester querying the oracle at the two sampled vectors and their sum
accepts is equal to the probability that the three-query BLR tester accepts $\ell$;
equivalently, both are the probability over independent uniform $x,y\in\mathbb{F}_2^n$
that $\ell(x)+\ell(y)=\ell(x+y)$.
\end{theorem}

\begin{theorem}[Acceptance probability of the matrix BLR tester]
\lean{FABL.quadraticPairMatrixBLRTester_acceptanceProbability}
\label{FABL.quadraticPairMatrixBLRTester_acceptanceProbability}
\leanok
\uses{FABL.F₂Cube, FABL.NonadaptiveFunctionTester.acceptanceProbability, FABL.blrFunctionTester, FABL.pmfExpectation, FABL.quadraticCharacterOracle, FABL.quadraticPairMatrixBLRTester, FABL.unflattenBinaryMatrix, FABL.uniformPMF, FABL.𝔽₂}
\difficulty{3}
For any natural number $n$, any function $\ell:\mathbb{F}_2^n\to\mathbb{F}_2$, and any
function $q:\mathbb{F}_2^{n\times n}\to\mathbb{F}_2$, the tester that independently
samples uniformly random binary matrices $A$ and $B$, queries the oracle at $A$, $B$,
and $A+B$, and accepts exactly when the first two answers sum to the third has the same
acceptance probability as the three-query BLR tester applied to the function on
$\mathbb{F}_2^{n^2}$ that sends a flattened matrix to its value under $q$.
\end{theorem}

\begin{definition}[Rank-one tester acceptance probability]
\lean{FABL.quadraticPairRankOneAcceptanceProbability}
\statementsource{Analysis-of-Boolean-Functions-by-Ryan-ODonnell}{pdf-ffbe9c7b3306-p198-b001}
\label{FABL.quadraticPairRankOneAcceptanceProbability}
\leanok
\uses{FABL.F₂Cube, FABL.binaryOuterProduct, FABL.unflattenBinaryMatrix, FABL.𝔽₂}
For $n\ge 0$, let $\ell:\mathbb F_2^n\to\mathbb F_2$ and $q:\mathbb F_2^{n\times
n}\to\mathbb F_2$. The rank-one tester's acceptance probability is the average, over
independent uniform $x,y\in\mathbb F_2^n$ and $z\in\mathbb F_2^{n^2}$, of the indicator
that $\ell(x)\ell(y)=q(Z)+q(Z+xy^{\mathsf T})$ in $\mathbb F_2$, where $Z$ is the matrix
obtained by reshaping $z$.
\end{definition}

\begin{theorem}[Tester acceptance probability formula]
\lean{FABL.quadraticPairRankOneTester_acceptanceProbability}
\label{FABL.quadraticPairRankOneTester_acceptanceProbability}
\leanok
\uses{FABL.F₂Cube, FABL.NonadaptiveFunctionTester.acceptanceProbability, FABL.NonadaptiveFunctionTester.transcript_apply, FABL.pmfExpectation, FABL.pmfExpectation_uniformPMF_eq_expect, FABL.quadraticCharacterOracle, FABL.quadraticPairRankOneAcceptanceProbability, FABL.quadraticPairRankOneTester, FABL.uniformPMF, FABL.𝔽₂}
\difficulty{4}
For a natural number $n$, let $\ell:\mathbb F_2^n\to\mathbb F_2$ and $q:\mathbb
F_2^{n\times n}\to\mathbb F_2$. The acceptance probability of the four-query tester for
the oracle that equals $\ell$ on vectors and $q$ on matrices is the uniform average over
independent uniform $x,y\in\mathbb F_2^n$ and $z\in\mathbb F_2^{n^2}$ of the indicator
that $\ell(x)\ell(y)=q(Z)+q(Z+xy^{\mathsf T})$, where $Z$ is the matrix obtained by
reshaping $z$ into an $n\times n$ matrix.
\end{theorem}

\begin{definition}[Quadratic character pair]
\lean{FABL.IsQuadraticCharacterPair}
\statementsource{Analysis-of-Boolean-Functions-by-Ryan-ODonnell}{
  pdf-ffbe9c7b3306-p198-b001,
  pdf-ffbe9c7b3306-p198-b003
}
\label{FABL.IsQuadraticCharacterPair}
\leanok
\uses{FABL.F₂Cube, FABL.binaryMatrixCharacter, FABL.f₂DotProduct, FABL.rankOneSelfMatrix, FABL.𝔽₂}
For $n\ge 0$, a pair of functions $\ell:\mathbb F_2^n\to\mathbb F_2$ and $q:\mathbb
F_2^{n\times n}\to\mathbb F_2$ is a quadratic character pair if there exists
$\gamma\in\mathbb F_2^n$ such that $\ell(x)=\sum_{i=1}^n\gamma_i x_i$ for every $x$, and
$q$ is the matrix character associated with the rank-one matrix $\gamma\gamma^{\mathsf
T}$.
\end{definition}

\begin{theorem}[Relative distance under matrix unflattening]
\lean{FABL.relativeHammingDist_unflattenBinaryMatrix}
\label{FABL.relativeHammingDist_unflattenBinaryMatrix}
\leanok
\uses{FABL.F₂Cube, FABL.binaryMatrixCubeEquiv, FABL.flattenBinaryMatrix, FABL.relativeHammingDist, FABL.unflattenBinaryMatrix, FABL.unflattenBinaryMatrix_flattenBinaryMatrix, FABL.uniformProbability, FABL.uniformProbability_ne_eq_relativeHammingDist, FABL.𝔽₂}
\difficulty{3}
Let $n$ be a nonnegative integer, and let $f,g$ be functions from the $n\times n$
matrices over the two-element field to that field. Regard each binary matrix as its
array of $n^2$ entries in a fixed coordinate order. The relative Hamming distance
between the resulting functions on binary arrays equals the relative Hamming distance
between $f$ and $g$ on matrices.
\end{theorem}

\begin{theorem}[Distance between linear Boolean forms]
\lean{FABL.relativeHammingDist_f₂DotProduct_eq}
\leanok
\uses{FABL.F₂Cube, FABL.f₂DotProduct, FABL.relativeHammingDist, FABL.uniformProbability, FABL.uniformProbability_f₂DotProduct_ne_zero, FABL.uniformProbability_ne_eq_relativeHammingDist}
\difficulty{4}
For any $a,b\in\mathbb F_2^n$, the relative Hamming distance between the functions
$x\mapsto a\cdot x$ and $x\mapsto b\cdot x$ on $\mathbb F_2^n$ is $0$ if $a=b$, and
is $1/2$ otherwise.
\end{theorem}

\begin{theorem}[Distance between binary matrix characters]
\lean{FABL.relativeHammingDist_binaryMatrixCharacter_eq}
\label{FABL.relativeHammingDist_binaryMatrixCharacter_eq}
\leanok
\uses{FABL.binaryMatrixCharacter, FABL.binaryMatrixCharacter_unflattenBinaryMatrix, FABL.binaryMatrixCubeEquiv, FABL.flattenBinaryMatrix, FABL.relativeHammingDist, FABL.relativeHammingDist_f₂DotProduct_eq, FABL.relativeHammingDist_unflattenBinaryMatrix, FABL.𝔽₂}
\difficulty{4}
Let $A$ and $B$ be square matrices of size $n$ over the field with two elements, and
associate to each matrix its character given by the sum of the entrywise products with
the input matrix. Their relative Hamming distance is $0$ if $A=B$ and $1/2$ otherwise.
\end{theorem}

\begin{theorem}[Triangle Inequality for Relative Hamming Distance]
\lean{FABL.relativeHammingDist_triangle_left}
\label{FABL.relativeHammingDist_triangle_left}
\leanok
\uses{FABL.relativeHammingDist}
\difficulty{3}
Let $\Omega$ be a finite nonempty set, let $\beta$ be a type with decidable equality,
and let $f,g,h:\Omega\to\beta$. Then the relative Hamming distance between $f$ and $g$
is at most the sum of the relative Hamming distances between $h$ and $f$ and between $h$
and $g$.
\end{theorem}

\begin{theorem}[Quadratic pair linear blr tester complete]
\lean{FABL.quadraticPairLinearBLRTester_complete}
\statementsource{Analysis-of-Boolean-Functions-by-Ryan-ODonnell}{pdf-ffbe9c7b3306-p197-b010}
\label{FABL.quadraticPairLinearBLRTester_complete}
\leanok
\uses{FABL.F₂Cube, FABL.NonadaptiveFunctionTester.acceptanceProbability, FABL.binaryMatrixCharacter, FABL.blrFunctionTester_complete, FABL.f₂DotProduct, FABL.isF₂Linear_iff_exists_dotProduct, FABL.quadraticCharacterOracle, FABL.quadraticPairLinearBLRTester, FABL.quadraticPairLinearBLRTester_acceptanceProbability, FABL.rankOneSelfMatrix}
\difficulty{3}
% dataset-pipeline: stub -- prose pending build_informal_statements.py
Statement pending regeneration.
\end{theorem}

\begin{theorem}[Completeness of the matrix linearity tester]
\lean{FABL.quadraticPairMatrixBLRTester_complete}
\label{FABL.quadraticPairMatrixBLRTester_complete}
\leanok
\uses{FABL.F₂Cube, FABL.NonadaptiveFunctionTester.acceptanceProbability, FABL.binaryMatrixCharacter, FABL.binaryMatrixCharacter_unflattenBinaryMatrix, FABL.blrFunctionTester_complete, FABL.flattenBinaryMatrix, FABL.f₂DotProduct, FABL.isF₂Linear_iff_exists_dotProduct, FABL.quadraticCharacterOracle, FABL.quadraticPairMatrixBLRTester, FABL.quadraticPairMatrixBLRTester_acceptanceProbability, FABL.rankOneSelfMatrix}
\difficulty{3}
For every natural number $n$ and every function $\gamma:\{1,\ldots,n\}\to\mathbb F_2$,
form the matrix $R_\gamma$ with entries $(R_\gamma)_{ij}=\gamma(i)\gamma(j)$. Choose
independent uniform matrices $A,B\in\mathbb F_2^{n\times n}$ and query the matrix oracle
$Q(C)=\sum_{i,j} (R_\gamma)_{ij}C_{ij}$ at $A$, $B$, and $A+B$; the tester accepts
exactly when the first two answers sum to the third. Its acceptance probability is $1$.
\end{theorem}

\begin{theorem}[Completeness of the rank-one tester]
\lean{FABL.quadraticPairRankOneTester_complete}
\label{FABL.quadraticPairRankOneTester_complete}
\leanok
\uses{FABL.F₂Cube, FABL.NonadaptiveFunctionTester.acceptanceProbability, FABL.binaryMatrixCharacter, FABL.binaryMatrixCharacter_add, FABL.binaryMatrixPairing, FABL.binaryMatrixPairing_rankOne, FABL.binaryOuterProduct, FABL.f₂DotProduct, FABL.matrixRankOnePairing_rankOneSelfMatrix, FABL.quadraticCharacterOracle, FABL.quadraticPairRankOneAcceptanceProbability, FABL.quadraticPairRankOneTester, FABL.quadraticPairRankOneTester_acceptanceProbability, FABL.rankOneSelfMatrix, FABL.unflattenBinaryMatrix}
\difficulty{4}
For every natural number $n$ and every $\gamma\in\mathbb F_2^n$, the four-query tester
accepts with probability $1$ the oracle that returns the dot product $\gamma\cdot x$ on
vectors $x$ and the matrix pairing $\sum_{i,j}\gamma_i\gamma_j A_{ij}$ on binary
matrices $A$.
\end{theorem}

\begin{theorem}[Perfect completeness for quadratic character pairs]
\lean{FABL.quadraticCharacterPairTester_complete}
\label{FABL.quadraticCharacterPairTester_complete}
\leanok
\uses{FABL.F₂Cube, FABL.IsQuadraticCharacterPair, FABL.NonadaptiveFunctionTester.acceptanceProbability, FABL.binaryMatrixCharacter, FABL.f₂DotProduct, FABL.quadraticCharacterOracle, FABL.quadraticCharacterPairTester, FABL.quadraticPairLinearBLRTester, FABL.quadraticPairLinearBLRTester_complete, FABL.quadraticPairMatrixBLRTester, FABL.quadraticPairMatrixBLRTester_complete, FABL.quadraticPairRankOneTester, FABL.quadraticPairRankOneTester_complete, FABL.rankOneSelfMatrix, FABL.uniformMixtureFunctionTester_acceptanceProbability, FABL.𝔽₂}
\difficulty{4}
Let $n\in\mathbb{N}$, let $\ell:\mathbb{F}_2^n\to\mathbb{F}_2$ and
$q:\mathbb{F}_2^{n\times n}\to\mathbb{F}_2$, and suppose there is
$\gamma\in\mathbb{F}_2^n$ such that $\ell(x)=\sum_i\gamma_i x_i$ for every $x$ and
$q(B)=\sum_{i,j}\gamma_i\gamma_jB_{ij}$ for every matrix $B$. Then the quadratic
character pair tester accepts the oracle formed from $\ell$ and $q$ with probability
exactly $1$. All sums and products are in $\mathbb{F}_2$.
\end{theorem}

\begin{theorem}[Matrix character mismatch under translation]
\lean{FABL.expect_matrixCharacter_mismatch_add}
\label{FABL.expect_matrixCharacter_mismatch_add}
\leanok
\uses{FABL.F₂Cube, FABL.binaryMatrixCharacter, FABL.flattenBinaryMatrix, FABL.relativeHammingDist, FABL.relativeHammingDist_unflattenBinaryMatrix, FABL.unflattenBinaryMatrix, FABL.unflattenBinaryMatrix_add, FABL.uniformProbability, FABL.uniformProbability_ne_eq_relativeHammingDist, FABL.𝔽₂}
\difficulty{4}
Let $n$ be a nonnegative integer, let $q$ be any function from $n\times n$ matrices over
$\mathbb F_2$ to $\mathbb F_2$, and let $A$ and $X$ be such matrices. For a uniformly
random binary matrix $B$, the probability that $q(B+X)$ differs from the matrix
character $C\mapsto\sum_{i,j}A_{ij}C_{ij}$ evaluated at $B+X$ equals the relative
Hamming distance between $q$ and that character.
\end{theorem}

\begin{theorem}[Rejection indicator bound]
\lean{FABL.rankOne_rejection_indicator_le_observed_add_query_errors}
\label{FABL.rankOne_rejection_indicator_le_observed_add_query_errors}
\leanok
\uses{FABL.𝔽₂}
\difficulty{5}
Let $\ell_x,\ell_y,\ell_x^0,\ell_y^0,q_z,q_z^x,q_z^0,q_z^{x,0}\in\mathbb F_2$. The
indicator of the event $\ell_x^0\ell_y^0\ne q_z^0+q_z^{x,0}$ is at most the indicator of
$\ell_x\ell_y\ne q_z+q_z^x$, plus the indicators of $\ell_x\ne\ell_x^0$,
$\ell_y\ne\ell_y^0$, $q_z\ne q_z^0$, and $q_z^x\ne q_z^{x,0}$.
\end{theorem}

\begin{theorem}[Rejection probability as a disagreement average]
\lean{FABL.quadraticPairRankOne_rejectionProbability_eq}
\label{FABL.quadraticPairRankOne_rejectionProbability_eq}
\leanok
\uses{FABL.F₂Cube, FABL.binaryOuterProduct, FABL.quadraticPairRankOneAcceptanceProbability, FABL.unflattenBinaryMatrix, FABL.𝔽₂}
\difficulty{4}
For every natural number $n$, every function $\ell:\mathbb F_2^n\to\mathbb F_2$, and
every function $q:\mathbb F_2^{n\times n}\to\mathbb F_2$, one minus the acceptance
probability defined by the uniform average over independent uniform $x,y\in\mathbb
F_2^n$ and $z\in\mathbb F_2^{n^2}$ equals the uniform average of the indicator that
$\ell(x)\ell(y)\ne q(Z)+q(Z+xy^{\mathsf T})$, with $Z$ the matrix obtained by reshaping
$z$.
\end{theorem}

\begin{theorem}[Complement of rank-one consistency probability]
\lean{FABL.one_sub_rankOneConsistencyProbability_eq}
\label{FABL.one_sub_rankOneConsistencyProbability_eq}
\leanok
\uses{FABL.F₂Cube, FABL.binaryMatrixCharacter, FABL.binaryOuterProduct, FABL.f₂DotProduct, FABL.matrixRankOnePairing, FABL.rankOneConsistencyProbability, FABL.uniformProbability, FABL.𝔽₂}
\difficulty{4}
Let $n$ be a natural number, let $\gamma\in\mathbb F_2^n$, and let $A$ be an $n\times n$
matrix over $\mathbb F_2$. The probability that, for independent uniform $x,y\in\mathbb
F_2^n$, the product of the dot products $\langle\gamma,x\rangle\langle\gamma,y\rangle$
differs from the matrix pairing $\sum_{i,j}A_{ij}x_i y_j$ is $1$ minus the probability
that these two quantities are equal.
\end{theorem}

\begin{theorem}[Consistency error bound from acceptance and distances]
\lean{FABL.one_sub_rankOneConsistencyProbability_le_observed_add_distances}
\label{FABL.one_sub_rankOneConsistencyProbability_le_observed_add_distances}
\leanok
\uses{FABL.F₂Cube, FABL.binaryMatrixCharacter, FABL.binaryMatrixCharacter_add, FABL.binaryOuterProduct, FABL.expect_matrixCharacter_mismatch_add, FABL.f₂DotProduct, FABL.one_sub_rankOneConsistencyProbability_eq, FABL.quadraticPairRankOneAcceptanceProbability, FABL.quadraticPairRankOne_rejectionProbability_eq, FABL.rankOneConsistencyProbability, FABL.rankOne_rejection_indicator_le_observed_add_query_errors, FABL.relativeHammingDist, FABL.relativeHammingDist_unflattenBinaryMatrix, FABL.unflattenBinaryMatrix, FABL.uniformProbability, FABL.uniformProbability_ne_eq_relativeHammingDist, FABL.𝔽₂}
\difficulty{5}
Let $n$ be a natural number, let $\ell:\mathbb F_2^n\to\mathbb F_2$ and $q:\mathbb
F_2^{n\times n}\to\mathbb F_2$, and choose $\gamma\in\mathbb F_2^n$ and an $n\times n$
matrix $A$ over $\mathbb F_2$. The rank-one consistency error for $\gamma,A$ is at most
the rejection probability of the test defined by $\ell,q$, plus twice the relative
Hamming distance from $\ell$ to $x\mapsto\langle\gamma,x\rangle$, plus twice the
relative Hamming distance from $q$ to the matrix character defined by $A$.
\end{theorem}

\begin{definition}[Quadratic character pair rejection rate]
\lean{FABL.quadraticCharacterPairRejectionRate}
\statementsource{Analysis-of-Boolean-Functions-by-Ryan-ODonnell}{
  pdf-ffbe9c7b3306-p198-b001,
  pdf-ffbe9c7b3306-p198-b003
}
\label{FABL.quadraticCharacterPairRejectionRate}
\leanok
For every $n\ge 0$, the quadratic character pair rejection rate is the real number
$1/100$.
\end{definition}

\begin{theorem}[Closeness from high tester acceptance]
\lean{FABL.quadraticCharacterPairTester_close_of_acceptanceProbability_ge}
\proofsource{Analysis-of-Boolean-Functions-by-Ryan-ODonnell}{
  pdf-ffbe9c7b3306-p198-b001,
  pdf-ffbe9c7b3306-p198-b003
}
\label{FABL.quadraticCharacterPairTester_close_of_acceptanceProbability_ge}
\leanok
\uses{FABL.F₂Cube, FABL.IsClose, FABL.IsLocalTester.close_of_acceptanceProbability_ge, FABL.NonadaptiveFunctionTester, FABL.NonadaptiveFunctionTester.acceptanceProbability, FABL.QuadraticCharacterAddress, FABL.binaryMatrixCharacter, FABL.binaryMatrixCharacter_unflattenBinaryMatrix, FABL.blrFunctionTester_isLocalTester, FABL.component_acceptanceProbability_ge_of_uniformMixture_ge, FABL.flattenBinaryMatrix, FABL.f₂DotProduct, FABL.isF₂Linear_iff_exists_dotProduct, FABL.one_sub_rankOneConsistencyProbability_le_observed_add_distances, FABL.quadraticCharacterOracle, FABL.quadraticCharacterPairRejectionRate, FABL.quadraticCharacterPairTester, FABL.quadraticPairLinearBLRTester, FABL.quadraticPairLinearBLRTester_acceptanceProbability, FABL.quadraticPairMatrixBLRTester, FABL.quadraticPairMatrixBLRTester_acceptanceProbability, FABL.quadraticPairRankOneAcceptanceProbability, FABL.quadraticPairRankOneTester, FABL.quadraticPairRankOneTester_acceptanceProbability, FABL.rankOneConsistencyProbability, FABL.rankOneConsistencyProbability_le_three_fourths, FABL.rankOneSelfMatrix, FABL.relativeHammingDist, FABL.relativeHammingDist_unflattenBinaryMatrix, FABL.unflattenBinaryMatrix, FABL.uniformMixtureFunctionTester, FABL.𝔽₂}
\difficulty{6}
Let $n\in\mathbb{N}$, let $0\le\varepsilon\le1$, and let
$\ell:\mathbb{F}_2^n\to\mathbb{F}_2$ and $q:\mathbb{F}_2^{n\times n}\to\mathbb{F}_2$. If
the quadratic character pair tester accepts the oracle that uses $\ell$ on vectors and
$q$ on matrices with probability at least $1-\varepsilon/100$, then there is
$\gamma\in\mathbb{F}_2^n$ such that the fraction of inputs on which $\ell(x)$ differs
from $\sum_i\gamma_i x_i$ is at most $\varepsilon$, and the fraction of matrices on
which $q(B)$ differs from $\sum_{i,j}\gamma_i\gamma_jB_{ij}$ is at most $\varepsilon$.
All sums and products in these formulas are in $\mathbb{F}_2$.
\end{theorem}

\begin{definition}[Homogeneous quadratic equation data]
\lean{FABL.HomogeneousQuadraticEquation}
\statementsource{Analysis-of-Boolean-Functions-by-Ryan-ODonnell}{pdf-ffbe9c7b3306-p198-b003}
\label{FABL.HomogeneousQuadraticEquation}
\leanok
\uses{FABL.𝔽₂}
For each $n\ge 0$, a homogeneous quadratic equation over $\mathbb F_2$ consists of an
$n\times n$ coefficient matrix over $\mathbb F_2$ and a right-hand side in $\mathbb
F_2$.
\end{definition}

\begin{definition}[Satisfaction of a homogeneous quadratic equation]
\lean{FABL.HomogeneousQuadraticEquation.SatisfiedBy}
\statementsource{Analysis-of-Boolean-Functions-by-Ryan-ODonnell}{pdf-ffbe9c7b3306-p198-b003}
\label{FABL.HomogeneousQuadraticEquation.SatisfiedBy}
\leanok
\uses{FABL.F₂Cube, FABL.HomogeneousQuadraticEquation, FABL.binaryMatrixPairing, FABL.rankOneSelfMatrix}
For $n\ge 0$, a homogeneous quadratic equation with coefficient matrix $A\in\mathbb
F_2^{n\times n}$ and right-hand side $b\in\mathbb F_2$ is satisfied by $w\in\mathbb
F_2^n$ when the matrix pairing of $A$ with $ww^{\mathsf T}$ equals $b$ in $\mathbb F_2$.
\end{definition}

\begin{definition}[Homogeneous quadratic system]
\lean{FABL.HomogeneousQuadraticSystem}
\statementsource{Analysis-of-Boolean-Functions-by-Ryan-ODonnell}{pdf-ffbe9c7b3306-p198-b003}
\label{FABL.HomogeneousQuadraticSystem}
\leanok
\uses{FABL.HomogeneousQuadraticEquation}
For a natural number $n$, a homogeneous quadratic system in $n$ variables over $\mathbb{F}_2$ is
a finite list of homogeneous quadratic equations in those variables.
\end{definition}

\begin{definition}[Solutions of a quadratic system]
\lean{FABL.SatisfiesHomogeneousQuadraticSystem}
\statementsource{Analysis-of-Boolean-Functions-by-Ryan-ODonnell}{pdf-ffbe9c7b3306-p198-b003}
\label{FABL.SatisfiesHomogeneousQuadraticSystem}
\leanok
\uses{FABL.F₂Cube, FABL.HomogeneousQuadraticEquation.SatisfiedBy, FABL.HomogeneousQuadraticSystem}
Let $n$ be a natural number, let $L$ be a finite list of homogeneous quadratic equations
in $n$ variables over the field with two elements, and let $w$ be an $n$-tuple over that
field. Then $w$ satisfies $L$ exactly when it satisfies every equation in the list.
\end{definition}

\begin{definition}[Solutions of the linearized system]
\lean{FABL.SatisfiesLinearizedQuadraticSystem}
\statementsource{Analysis-of-Boolean-Functions-by-Ryan-ODonnell}{
  pdf-ffbe9c7b3306-p198-b001,
  pdf-ffbe9c7b3306-p198-b003
}
\label{FABL.SatisfiesLinearizedQuadraticSystem}
\leanok
\uses{FABL.F₂Cube, FABL.HomogeneousQuadraticSystem, FABL.flattenBinaryMatrix, FABL.f₂DotProduct}
Let $n$ be a natural number, let $L$ be a finite list of homogeneous quadratic equations
in $n$ variables over the field with two elements, and let $a$ be a vector of length
$n^2$ over that field. Then $a$ satisfies the linearized system exactly when, for every
equation in $L$, its coefficient matrix flattened into a vector has dot product with $a$
equal to that equation’s right-hand side.
\end{definition}

\begin{definition}[Direction subspace of a quadratic system]
\lean{FABL.homogeneousQuadraticSystemDirection}
\statementsource{Analysis-of-Boolean-Functions-by-Ryan-ODonnell}{pdf-ffbe9c7b3306-p198-b003}
\label{FABL.homogeneousQuadraticSystemDirection}
\leanok
\uses{FABL.F₂Cube, FABL.HomogeneousQuadraticSystem, FABL.flattenBinaryMatrix, FABL.f₂DotProduct, FABL.𝔽₂}
Let $n$ be a natural number and let $L$ be a finite list of homogeneous quadratic
equations in $n$ variables over the field with two elements. The associated direction
subspace consists of the vectors $a$ of length $n^2$ for which the dot product of $a$
with every flattened coefficient matrix from $L$ is zero.
\end{definition}

\begin{theorem}[Rank-One Linearization Equivalence]
\lean{FABL.satisfiesLinearized_rankOne_iff}
\label{FABL.satisfiesLinearized_rankOne_iff}
\leanok
\uses{FABL.F₂Cube, FABL.HomogeneousQuadraticSystem, FABL.SatisfiesHomogeneousQuadraticSystem, FABL.SatisfiesLinearizedQuadraticSystem, FABL.flattenBinaryMatrix, FABL.f₂DotProduct_flattenBinaryMatrix, FABL.rankOneSelfMatrix}
\difficulty{3}
Let $L$ be a list of homogeneous quadratic equations over $\mathbb{F}_2$ in $n$
variables, and let $w\in\mathbb{F}_2^n$. The rank-one matrix formed from $w$, when
flattened into a vector, satisfies the linearized system associated with $L$ if and only
if $w$ satisfies every equation in $L$.
\end{theorem}

\begin{theorem}[Linearized Solutions as an Affine Subspace]
\lean{FABL.satisfiesLinearized_iff_mem_binaryAffineSubspace}
\label{FABL.satisfiesLinearized_iff_mem_binaryAffineSubspace}
\leanok
\uses{FABL.F₂Cube, FABL.HomogeneousQuadraticSystem, FABL.SatisfiesLinearizedQuadraticSystem, FABL.binaryAffineSubspace, FABL.f₂DotProduct, FABL.homogeneousQuadraticSystemDirection, FABL.mem_binaryAffineSubspace_iff_add_mem, FABL.𝔽₂}
\difficulty{4}
Let $L$ be a list of homogeneous quadratic equations in $n$ variables over
$\mathbb{F}_2$, and suppose that $a_0\in\mathbb{F}_2^{n^2}$ satisfies the associated
linearized system. For every $a\in\mathbb{F}_2^{n^2}$, $a$ satisfies that system if and
only if it belongs to the affine subspace through $a_0$ whose direction consists of the
solutions to the corresponding homogeneous linear system.
\end{theorem}

\begin{theorem}[Zero-dimensional quadratic equation]
\lean{FABL.homogeneousQuadraticEquation_satisfiedBy_zero_iff}
\label{FABL.homogeneousQuadraticEquation_satisfiedBy_zero_iff}
\leanok
\uses{FABL.F₂Cube, FABL.HomogeneousQuadraticEquation, FABL.HomogeneousQuadraticEquation.SatisfiedBy, FABL.binaryMatrixPairing}
\difficulty{1}
For the zero-dimensional binary cube, a homogeneous quadratic equation is satisfied by
its unique vector if and only if its right-hand side is zero.
\end{theorem}

\begin{theorem}[Zero-Dimensional Quadratic Systems]
\lean{FABL.satisfiesHomogeneousQuadraticSystem_zero_iff}
\label{FABL.satisfiesHomogeneousQuadraticSystem_zero_iff}
\leanok
\uses{FABL.F₂Cube, FABL.HomogeneousQuadraticSystem, FABL.SatisfiesHomogeneousQuadraticSystem, FABL.homogeneousQuadraticEquation_satisfiedBy_zero_iff}
\difficulty{4}
Let $L$ be a finite list of homogeneous quadratic equations over $\mathbb{F}_2$ in zero
variables, and let $w$ be the unique assignment on the empty set of variables. Then $w$
satisfies every equation in $L$ if and only if every equation in $L$ has right-hand side
$0$.
\end{theorem}

\begin{theorem}[Zero-dimensional system dichotomy]
\lean{FABL.homogeneousQuadraticSystem_zero_universal_or_empty}
\label{FABL.homogeneousQuadraticSystem_zero_universal_or_empty}
\leanok
\uses{FABL.HomogeneousQuadraticSystem, FABL.SatisfiesHomogeneousQuadraticSystem, FABL.satisfiesHomogeneousQuadraticSystem_zero_iff}
\difficulty{4}
Let $L$ be a finite system of homogeneous quadratic equations in zero variables over
$\mathbb{F}_2$. Either every zero-dimensional Boolean word satisfies every equation in
$L$, or no zero-dimensional Boolean word satisfies every equation in $L$.
\end{theorem}

\begin{definition}[Boolean solutions of a quadratic system]
\lean{FABL.IsBoolHomogeneousQuadraticSystem}
\statementsource{Analysis-of-Boolean-Functions-by-Ryan-ODonnell}{pdf-ffbe9c7b3306-p198-b003}
\label{FABL.IsBoolHomogeneousQuadraticSystem}
\leanok
\uses{FABL.HomogeneousQuadraticSystem, FABL.SatisfiesHomogeneousQuadraticSystem, FABL.pcppBoolCubeEquiv}
Let $n$ be a natural number, let $L$ be a finite list of homogeneous quadratic equations
in $n$ variables over the field with two elements, and let $w$ be a Boolean vector of
length $n$. Then $w$ satisfies the system when its image under the Boolean-to-field
coordinate map satisfies every equation in $L$.
\end{definition}

\begin{definition}[Tester with binary outputs from sign outputs]
\lean{FABL.f₂FunctionTesterOfSign}
\leanok
\uses{FABL.NonadaptiveFunctionTester, FABL.Sign, FABL.signEncode, FABL.𝔽₂}
Given a finite domain $\Omega$, a number of queries $r$, and a nonadaptive tester whose
outputs are signs, define a tester with outputs in $\mathbb{F}_2$. It uses the same
random seeds, seed distribution, and queries; on each answered transcript it converts
every sign value to its binary encoding before applying the original acceptance rule.
\end{definition}

\begin{theorem}[Acceptance probability under binary encoding]
\lean{FABL.f₂FunctionTesterOfSign_acceptanceProbability}
\leanok
\uses{FABL.NonadaptiveFunctionTester, FABL.NonadaptiveFunctionTester.acceptanceProbability, FABL.NonadaptiveFunctionTester.transcript_apply, FABL.Sign, FABL.f₂FunctionTesterOfSign, FABL.signEncode, FABL.𝔽₂}
\difficulty{4}
Let $\Omega$ be a finite set, let $r$ be a natural number, let $T$ be a nonadaptive
function tester from $\Omega$ to the sign group $\{\pm1\}$ with $r$ transcript
positions, and let $f:\Omega\to\mathbb F_2$. The acceptance probability of the tester
obtained by converting $T$'s answers from signs to bits on input $f$ equals the
acceptance probability of $T$ on the sign-valued function obtained by encoding each
value of $f$.
\end{theorem}

\begin{theorem}[Sign Encoding Preserves Relative Distance]
\lean{FABL.relativeHammingDist_signEncode}
\proofsource{Analysis-of-Boolean-Functions-by-Ryan-ODonnell}{
  pdf-ffbe9c7b3306-p026-b003,
  pdf-ffbe9c7b3306-p026-b004
}
\label{FABL.relativeHammingDist_signEncode}
\leanok
\uses{FABL.binarySignEquiv, FABL.relativeHammingDist, FABL.signEncode, FABL.𝔽₂}
\difficulty{3}
Let $\Omega$ be a finite nonempty set and let $f,g:\Omega\to\mathbb{F}_2$. Encoding each
value by its sign preserves relative Hamming distance: the relative Hamming distance
between the encoded functions equals that between $f$ and $g$.
\end{theorem}

\begin{definition}[Always rejecting tester on the binary cube]
\lean{FABL.alwaysRejectF₂CubeFunctionTester}
\leanok
\uses{FABL.F₂Cube, FABL.NonadaptiveFunctionTester, FABL.uniformPMF, FABL.𝔽₂}
For natural numbers $d$ and $r$, define a nonadaptive tester on the $d$-dimensional
binary cube with $r$ query positions. It has a single random seed, makes no queries, and
rejects every transcript.
\end{definition}

\begin{theorem}[Acceptance probability of an always rejecting tester]
\lean{FABL.alwaysRejectF₂CubeFunctionTester_acceptanceProbability}
\leanok
\uses{FABL.F₂Cube, FABL.NonadaptiveFunctionTester.acceptanceProbability, FABL.alwaysRejectF₂CubeFunctionTester, FABL.pmfExpectation, FABL.pmfExpectation_const, FABL.uniformPMF, FABL.𝔽₂}
\difficulty{3}
For any nonnegative integers $d,r$ and any function $f:\{0,1\}^d\to\mathbb F_2$, the
tester that always rejects has acceptance probability $0$ on $f$.
\end{theorem}

\begin{definition}[Character tester for the linearized system]
\lean{FABL.linearizedQuadraticSystemCharacterTester}
\statementsource{Analysis-of-Boolean-Functions-by-Ryan-ODonnell}{
  pdf-ffbe9c7b3306-p198-b001,
  pdf-ffbe9c7b3306-p198-b003
}
\label{FABL.linearizedQuadraticSystemCharacterTester}
\leanok
\uses{FABL.F₂Cube, FABL.HomogeneousQuadraticSystem, FABL.NonadaptiveFunctionTester, FABL.SatisfiesLinearizedQuadraticSystem, FABL.affineCharacterFunctionTester, FABL.alwaysRejectF₂CubeFunctionTester, FABL.f₂FunctionTesterOfSign, FABL.homogeneousQuadraticSystemDirection, FABL.𝔽₂}
Let $n$ be a natural number and let $L$ be a finite list of homogeneous quadratic
equations in $n$ variables over the field with two elements. The associated nonadaptive
tester has three query positions, queries functions on vectors of length $n^2$ over that
field with values in the field, and uses a finite seed space with a probability
distribution. It is chosen by whether the linearized system has a solution: when one
exists, the tester is built from the affine-character tester for the direction subspace
using such a solution; otherwise, it always rejects.
\end{definition}

\begin{theorem}[Completeness of the character tester]
\lean{FABL.linearizedQuadraticSystemCharacterTester_complete}
\label{FABL.linearizedQuadraticSystemCharacterTester_complete}
\leanok
\uses{FABL.F₂Cube, FABL.HomogeneousQuadraticSystem, FABL.NonadaptiveFunctionTester.acceptanceProbability, FABL.SatisfiesLinearizedQuadraticSystem, FABL.affineCharacterFunctionTester_isLocalTester, FABL.f₂DotProduct, FABL.f₂FunctionTesterOfSign_acceptanceProbability, FABL.homogeneousQuadraticSystemDirection, FABL.linearizedQuadraticSystemCharacterTester, FABL.satisfiesLinearized_iff_mem_binaryAffineSubspace}
\difficulty{4}
Let $L$ be a finite system of homogeneous quadratic equations over $\mathbb F_2$, and
let $a\in\mathbb F_2^{n^2}$ satisfy every equation of $L$. Then the linearized-system
character tester for $L$ accepts the function $x\mapsto a\cdot x$ with probability
exactly $1$.
\end{theorem}

\begin{theorem}[Closeness from high tester acceptance]
\lean{FABL.linearizedQuadraticSystemCharacterTester_close_of_acceptanceProbability_ge}
\label{FABL.linearizedQuadraticSystemCharacterTester_close_of_acceptanceProbability_ge}
\leanok
\uses{FABL.F₂Cube, FABL.HomogeneousQuadraticSystem, FABL.IsClose, FABL.IsLocalTester.close_of_acceptanceProbability_ge, FABL.NonadaptiveFunctionTester.acceptanceProbability, FABL.SatisfiesLinearizedQuadraticSystem, FABL.affineCharacterFunctionTester, FABL.affineCharacterFunctionTester_isLocalTester, FABL.alwaysRejectF₂CubeFunctionTester_acceptanceProbability, FABL.f₂DotProduct, FABL.f₂FunctionTesterOfSign_acceptanceProbability, FABL.homogeneousQuadraticSystemDirection, FABL.linearizedQuadraticSystemCharacterTester, FABL.relativeHammingDist, FABL.relativeHammingDist_signEncode, FABL.satisfiesLinearized_iff_mem_binaryAffineSubspace, FABL.signEncode, FABL.𝔽₂}
\difficulty{5}
Let $L$ be a finite system of homogeneous quadratic equations over $\mathbb F_2$, let
$0\leq\varepsilon\leq1$, and let $f:\mathbb F_2^{n^2}\to\mathbb F_2$. If the
linearized-system character tester for $L$ accepts $f$ with probability at least
$1-\varepsilon/2$, then there is a vector $a\in\mathbb F_2^{n^2}$ satisfying every
equation of $L$ such that the relative Hamming distance from $f$ to the linear function
$x\mapsto a\cdot x$ is at most $\varepsilon$.
\end{theorem}

\begin{definition}[Index for the first block]
\lean{FABL.quadraticLinearProofIndex}
\statementsource{Analysis-of-Boolean-Functions-by-Ryan-ODonnell}{
  pdf-ffbe9c7b3306-p198-b003,
  pdf-ffbe9c7b3306-p198-b004
}
\label{FABL.quadraticLinearProofIndex}
\leanok
Let $n$ be a natural number. The first-block index map sends each index in a block of
size $2^n$ to the corresponding index in the first summand of a combined index set of
size $2^n+2^{n^2}$.
\end{definition}

\begin{definition}[Index for the matrix block]
\lean{FABL.quadraticMatrixProofIndex}
\statementsource{Analysis-of-Boolean-Functions-by-Ryan-ODonnell}{
  pdf-ffbe9c7b3306-p198-b003,
  pdf-ffbe9c7b3306-p198-b004
}
\label{FABL.quadraticMatrixProofIndex}
\leanok
Let $n$ be a natural number. The matrix-block index map sends each index in a block of
size $2^{n^2}$ to the corresponding index in the second summand of a combined index set
of size $2^n+2^{n^2}$.
\end{definition}

\begin{definition}[Linear proof table restriction]
\lean{FABL.quadraticLinearProofTable}
\statementsource{Analysis-of-Boolean-Functions-by-Ryan-ODonnell}{
  pdf-ffbe9c7b3306-p197-b012,
  pdf-ffbe9c7b3306-p198-b003,
  pdf-ffbe9c7b3306-p198-b004
}
\label{FABL.quadraticLinearProofTable}
\leanok
\uses{FABL.quadraticLinearProofIndex}
Let $n$ be a natural number and let $P$ be a Boolean table indexed by the disjoint union
of sets of sizes $2^n$ and $2^{n^2}$. Its restriction to the first summand is a Boolean
table on a set of size $2^n$, obtained by evaluating $P$ at the canonical inclusion of
each index into that summand.
\end{definition}

\begin{definition}[Matrix proof table restriction]
\lean{FABL.quadraticMatrixProofTable}
\statementsource{Analysis-of-Boolean-Functions-by-Ryan-ODonnell}{
  pdf-ffbe9c7b3306-p198-b003,
  pdf-ffbe9c7b3306-p198-b004
}
\label{FABL.quadraticMatrixProofTable}
\leanok
\uses{FABL.quadraticMatrixProofIndex}
Let $n$ be a natural number and let $P$ be a Boolean table indexed by the disjoint union
of sets of sizes $2^n$ and $2^{n^2}$. Its restriction to the second summand is a Boolean
table on a set of size $2^{n^2}$, obtained by evaluating $P$ at the canonical inclusion
of each index into that summand.
\end{definition}

\begin{definition}[Linear proof function]
\lean{FABL.quadraticLinearProofFunction}
\statementsource{Analysis-of-Boolean-Functions-by-Ryan-ODonnell}{
  pdf-ffbe9c7b3306-p197-b012,
  pdf-ffbe9c7b3306-p198-b003,
  pdf-ffbe9c7b3306-p198-b004
}
\label{FABL.quadraticLinearProofFunction}
\leanok
\uses{FABL.F₂Cube, FABL.f₂CubeIndexEquiv, FABL.pcppBoolF₂Equiv, FABL.quadraticLinearProofTable, FABL.𝔽₂}
Let $n$ be a natural number and let $P$ be a Boolean table indexed by the disjoint union
of sets of sizes $2^n$ and $2^{n^2}$. The associated function on the binary cube of
dimension $n$ evaluates $P$ at the index of its input in the first summand, then
interprets the resulting Boolean value as an element of the field with two elements.
\end{definition}

\begin{definition}[Matrix proof function]
\lean{FABL.quadraticMatrixProofFunction}
\statementsource{Analysis-of-Boolean-Functions-by-Ryan-ODonnell}{
  pdf-ffbe9c7b3306-p198-b003,
  pdf-ffbe9c7b3306-p198-b004
}
\label{FABL.quadraticMatrixProofFunction}
\leanok
\uses{FABL.flattenBinaryMatrix, FABL.f₂CubeIndexEquiv, FABL.pcppBoolF₂Equiv, FABL.quadraticMatrixProofTable, FABL.𝔽₂}
Let $n$ be a natural number and let $P$ be a Boolean table indexed by the disjoint union
of sets of sizes $2^n$ and $2^{n^2}$. The associated function on $n\times n$ matrices
over the field with two elements evaluates $P$ at the index of the flattened matrix in
the second summand, then interprets the resulting Boolean value as an element of the
field with two elements.
\end{definition}

\begin{definition}[Encoding a binary quadratic form]
\lean{FABL.encodeHomogeneousQuadraticProof}
\statementsource{Analysis-of-Boolean-Functions-by-Ryan-ODonnell}{
  pdf-ffbe9c7b3306-p198-b003,
  pdf-ffbe9c7b3306-p198-b004
}
\label{FABL.encodeHomogeneousQuadraticProof}
\leanok
\uses{FABL.F₂Cube, FABL.flattenBinaryMatrix, FABL.f₂CubeIndexEquiv, FABL.f₂DotProduct, FABL.pcppBoolF₂Equiv, FABL.rankOneSelfMatrix}
Let $n$ be a natural number and let $w\in\mathbb{F}_2^n$. Define a Boolean-valued
function on a disjoint union of an index set of size $2^n$ and one of size $2^{n^2}$ as
follows: at an index corresponding to $x\in\mathbb{F}_2^n$, its value is the Boolean
encoding of $\sum_{i=1}^n w_i x_i$ in $\mathbb{F}_2$; at an index corresponding to a
binary matrix $A=(A_{ij})_{1\le i,j\le n}$, its value is the Boolean encoding of
$\sum_{i,j=1}^n w_iw_jA_{ij}$ in $\mathbb{F}_2$.
\end{definition}

\begin{theorem}[Encoded proof recovers the linear function]
\lean{FABL.quadraticLinearProofFunction_encodeHomogeneousQuadraticProof}
\label{FABL.quadraticLinearProofFunction_encodeHomogeneousQuadraticProof}
\leanok
\uses{FABL.F₂Cube, FABL.encodeHomogeneousQuadraticProof, FABL.f₂DotProduct, FABL.quadraticLinearProofFunction, FABL.quadraticLinearProofIndex, FABL.quadraticLinearProofTable}
\difficulty{1}
For every binary vector $w\in\bbr^n$ over $\mathbb{F}_2$, the linear function extracted
from its encoded homogeneous quadratic proof is the dot product with $w$: at every
$x\in\mathbb{F}_2^n$, its value is $w\cdot x$.
\end{theorem}

\begin{theorem}[Quadratic matrix character encoding]
\lean{FABL.quadraticMatrixProofFunction_encodeHomogeneousQuadraticProof}
\proofsource{Analysis-of-Boolean-Functions-by-Ryan-ODonnell}{pdf-ffbe9c7b3306-p198-b003}
\label{FABL.quadraticMatrixProofFunction_encodeHomogeneousQuadraticProof}
\leanok
\uses{FABL.F₂Cube, FABL.binaryMatrixCharacter, FABL.encodeHomogeneousQuadraticProof, FABL.f₂DotProduct_flattenBinaryMatrix, FABL.quadraticMatrixProofFunction, FABL.quadraticMatrixProofIndex, FABL.quadraticMatrixProofTable, FABL.rankOneSelfMatrix}
\difficulty{3}
For every nonnegative integer $n$ and every vector $w\in\mathbb{F}_2^n$, the function on
$n\times n$ matrices over $\mathbb{F}_2$ obtained by decoding the homogeneous quadratic
proof encoding of $w$ is the matrix character given by pairing with the rank-one matrix
whose $(i,j)$ entry is $w_iw_j$.
\end{theorem}

\begin{definition}[Tester with reindexed query domain]
\lean{FABL.reindexDomainFunctionTester}
\label{FABL.reindexDomainFunctionTester}
\leanok
\uses{FABL.NonadaptiveFunctionTester}
Given types $\Omega$ and $\Omega'$, a codomain $\beta$, a natural number $r$, a map $e :
\Omega \to \Omega'$, and a nonadaptive tester from $\Omega$ to $\beta$, define a tester
from $\Omega'$ to $\beta$. It keeps the same random seeds, seed distribution, and
acceptance rule, and maps each queried address through $e$, preserving unqueried
positions.
\end{definition}

\begin{theorem}[Acceptance Probability Under Domain Reindexing]
\lean{FABL.reindexDomainFunctionTester_acceptanceProbability}
\label{FABL.reindexDomainFunctionTester_acceptanceProbability}
\leanok
\uses{FABL.NonadaptiveFunctionTester, FABL.NonadaptiveFunctionTester.acceptanceProbability, FABL.NonadaptiveFunctionTester.transcript_apply, FABL.reindexDomainFunctionTester}
\difficulty{4}
Let $\Omega,\Omega',\beta$ be types, let $r$ be a natural number, let
$e\colon\Omega\to\Omega'$ be any map, let $T$ be a nonadaptive tester on domain $\Omega$
with output type $\beta$ and $r$ query positions, and let $f\colon\Omega'\to\beta$. The
acceptance probability of the tester obtained by mapping each query through $e$, on
input function $f$, equals the acceptance probability of $T$ on the composite function
$f\circ e$.
\end{theorem}

\begin{definition}[Boolean-output tester from a binary-field tester]
\lean{FABL.boolFunctionTesterOfF₂}
\leanok
\uses{FABL.NonadaptiveFunctionTester, FABL.pcppBoolF₂Equiv, FABL.𝔽₂}
Given a type $\Omega$, a natural number $r$, and a nonadaptive tester from $\Omega$ to
$\mathbb{F}_2$, define a tester from $\Omega$ to the Booleans. It keeps the same random
seeds, seed distribution, and queries, converting each answered binary-field value to
its corresponding Boolean value before applying the original acceptance rule.
\end{definition}

\begin{theorem}[Acceptance under Boolean encoding]
\lean{FABL.boolFunctionTesterOfF₂_acceptanceProbability}
\leanok
\uses{FABL.NonadaptiveFunctionTester, FABL.NonadaptiveFunctionTester.acceptanceProbability, FABL.NonadaptiveFunctionTester.transcript_apply, FABL.boolFunctionTesterOfF₂, FABL.pcppBoolF₂Equiv, FABL.𝔽₂}
\difficulty{4}
Let $\Omega$ be a set, let $r$ be a natural number, let $T$ be a nonadaptive tester from
functions $\Omega\to\mathbb{F}_2$ to Boolean acceptance, and let
$f:\Omega\to\{\mathrm{false},\mathrm{true}\}$. The tester obtained by encoding Boolean
transcript answers as elements of $\mathbb{F}_2$ has, on $f$, the same acceptance
probability as $T$ on the function obtained by encoding each value of $f$ in
$\mathbb{F}_2$.
\end{theorem}

\begin{definition}[Padding a three-query tester]
\lean{FABL.padThreeQueryFunctionTester}
\label{FABL.padThreeQueryFunctionTester}
\leanok
\uses{FABL.NonadaptiveFunctionTester}
Given types $\Omega$ and $\beta$, and a nonadaptive tester from $\Omega$ to $\beta$ with
three query positions, define a tester with four query positions. It preserves the
random seeds and their distribution, uses the original three queries followed by no
query, and applies the original acceptance rule to the first three transcript entries.
\end{definition}

\begin{theorem}[Acceptance probability after adding a query]
\lean{FABL.padThreeQueryFunctionTester_acceptanceProbability}
\label{FABL.padThreeQueryFunctionTester_acceptanceProbability}
\leanok
\uses{FABL.NonadaptiveFunctionTester, FABL.NonadaptiveFunctionTester.acceptanceProbability, FABL.NonadaptiveFunctionTester.transcript_apply, FABL.padThreeQueryFunctionTester}
\difficulty{4}
Let $T$ be a nonadaptive function tester with three query positions, and let $f$ be a
function from its domain to its codomain. The tester obtained by padding $T$ with a
fourth, unused query position has the same acceptance probability on $f$ as $T$.
\end{theorem}

\begin{definition}[Quadratic pair proof address]
\lean{FABL.quadraticPairProofAddress}
\statementsource{Analysis-of-Boolean-Functions-by-Ryan-ODonnell}{
  pdf-ffbe9c7b3306-p198-b003,
  pdf-ffbe9c7b3306-p198-b004
}
\label{FABL.quadraticPairProofAddress}
\leanok
\uses{FABL.QuadraticCharacterAddress, FABL.flattenBinaryMatrix, FABL.f₂CubeIndexEquiv, FABL.quadraticLinearProofIndex, FABL.quadraticMatrixProofIndex}
Let $n$ be a natural number. This map sends a binary vector of length $n$ to the right
summand, using its canonical binary-cube index in the first block of size $2^n$; it
sends an $n\times n$ matrix over the field with two elements to the right summand, using
the canonical index of its flattened entries in the second block of size $2^{n^2}$. The
target is the disjoint union of a set of size $n$ and the proof-address set of size
$2^n+2^{n^2}$.
\end{definition}

\begin{definition}[Quadratic pair function tester]
\lean{FABL.quadraticPairProofFunctionTester}
\statementsource{Analysis-of-Boolean-Functions-by-Ryan-ODonnell}{
  pdf-ffbe9c7b3306-p198-b003,
  pdf-ffbe9c7b3306-p198-b004
}
\label{FABL.quadraticPairProofFunctionTester}
\leanok
\uses{FABL.NonadaptiveFunctionTester, FABL.boolFunctionTesterOfF₂, FABL.quadraticCharacterPairTester, FABL.quadraticPairProofAddress, FABL.reindexDomainFunctionTester}
For each natural number $n$, there is a randomized nonadaptive tester with four query
positions, domain the disjoint union of a set of size $n$ and the proof-address set of
size $2^n+2^{n^2}$, and Boolean codomain. It is obtained by reindexing the quadratic
character pair tester and converting its field-with-two-elements answers to Boolean
answers.
\end{definition}

\begin{theorem}[Acceptance probability under proof reindexing]
\lean{FABL.quadraticPairProofFunctionTester_acceptanceProbability}
\label{FABL.quadraticPairProofFunctionTester_acceptanceProbability}
\leanok
\uses{FABL.NonadaptiveFunctionTester.acceptanceProbability, FABL.boolFunctionTesterOfF₂_acceptanceProbability, FABL.inputProofWord, FABL.quadraticCharacterOracle, FABL.quadraticCharacterPairTester, FABL.quadraticLinearProofFunction, FABL.quadraticLinearProofTable, FABL.quadraticMatrixProofFunction, FABL.quadraticMatrixProofTable, FABL.quadraticPairProofAddress, FABL.quadraticPairProofFunctionTester, FABL.reindexDomainFunctionTester_acceptanceProbability}
\difficulty{4}
For every natural number $n$, every Boolean word
$w\colon\operatorname{Fin}(n)\to\{\mathrm{false},\mathrm{true}\}$, and every Boolean
proof table $P\colon\operatorname{Fin}(2^n+2^{n^2})\to\{\mathrm{false},\mathrm{true}\}$,
the proof-function tester's acceptance probability on the word that equals $w$ on input
coordinates and $P$ on proof coordinates is equal to the quadratic-character pair
tester's acceptance probability on the oracle whose linear and matrix components are
extracted from $P$.
\end{theorem}

\begin{definition}[Matrix cube proof address]
\lean{FABL.quadraticMatrixCubeProofAddress}
\statementsource{Analysis-of-Boolean-Functions-by-Ryan-ODonnell}{
  pdf-ffbe9c7b3306-p198-b003,
  pdf-ffbe9c7b3306-p198-b004
}
\label{FABL.quadraticMatrixCubeProofAddress}
\leanok
\uses{FABL.F₂Cube, FABL.f₂CubeIndexEquiv, FABL.quadraticMatrixProofIndex}
Let $n$ be a natural number and let $x$ be a vector of length $n^2$ over the field with
two elements. This map sends $x$ to the right summand of the disjoint union of a set of
size $n$ and a proof-address set of size $2^n+2^{n^2}$, at the position given by the
canonical binary-cube index of $x$ in the block of size $2^{n^2}$.
\end{definition}

\begin{definition}[Quadratic system proof tester]
\lean{FABL.quadraticSystemProofFunctionTester}
\statementsource{Analysis-of-Boolean-Functions-by-Ryan-ODonnell}{
  pdf-ffbe9c7b3306-p198-b003,
  pdf-ffbe9c7b3306-p198-b004
}
\label{FABL.quadraticSystemProofFunctionTester}
\leanok
\uses{FABL.HomogeneousQuadraticSystem, FABL.NonadaptiveFunctionTester, FABL.boolFunctionTesterOfF₂, FABL.linearizedQuadraticSystemCharacterTester, FABL.padThreeQueryFunctionTester, FABL.quadraticMatrixCubeProofAddress, FABL.reindexDomainFunctionTester}
Let $n$ be a natural number and let $L$ be a finite list of equations in $n$ variables
over the field with two elements, each specified by a coefficient matrix and a
right-hand side. The associated tester uses four query positions, takes addresses from
the disjoint union of the $n$ input positions and $2^n+2^{n^2}$ proof positions, and has
Boolean answers; it is the tester constructed to check the linearized quadratic system.
\end{definition}

\begin{theorem}[Proof tester acceptance probability]
\lean{FABL.quadraticSystemProofFunctionTester_acceptanceProbability}
\label{FABL.quadraticSystemProofFunctionTester_acceptanceProbability}
\leanok
\uses{FABL.HomogeneousQuadraticSystem, FABL.NonadaptiveFunctionTester.acceptanceProbability, FABL.boolFunctionTesterOfF₂_acceptanceProbability, FABL.inputProofWord, FABL.linearizedQuadraticSystemCharacterTester, FABL.padThreeQueryFunctionTester_acceptanceProbability, FABL.quadraticMatrixCubeProofAddress, FABL.quadraticMatrixProofFunction, FABL.quadraticMatrixProofTable, FABL.quadraticSystemProofFunctionTester, FABL.reindexDomainFunctionTester_acceptanceProbability, FABL.unflattenBinaryMatrix}
\difficulty{3}
Let $n$ be a natural number, let $L$ be a list of homogeneous quadratic equations in $n$
variables over $\mathbb F_2$, let
$w:\{0,\ldots,n-1\}\to\{\mathrm{false},\mathrm{true}\}$, and let
$p:\{0,\ldots,2^n+2^{n^2}-1\}\to\{\mathrm{false},\mathrm{true}\}$. The acceptance
probability of the four-query proof tester for $L$ on the combined word consisting of
$w$ and $p$ equals the acceptance probability of the three-query character tester for
$L$ on the function that maps each binary matrix $A$ to the value of $p$ at the
proof-table position indexed by the flattening of $A$.
\end{theorem}

\begin{definition}[Quadratic input consistency addresses]
\lean{FABL.quadraticInputConsistencyAddress}
\statementsource{Analysis-of-Boolean-Functions-by-Ryan-ODonnell}{pdf-ffbe9c7b3306-p198-b004}
\label{FABL.quadraticInputConsistencyAddress}
\leanok
\uses{FABL.quadraticLinearProofIndex}
For a natural number $n$, map the disjoint union of the $n$ input addresses and the
$2^n$ linear proof addresses into the disjoint union of the $n$ input addresses and the
$2^n+2^{n^2}$ proof addresses. Each input address is unchanged, and each linear proof
address is sent to its corresponding initial proof address among the first $2^n$
positions.
\end{definition}

\begin{definition}[Quadratic input consistency tester]
\lean{FABL.quadraticInputConsistencyFunctionTester}
\statementsource{Analysis-of-Boolean-Functions-by-Ryan-ODonnell}{pdf-ffbe9c7b3306-p198-b004}
\label{FABL.quadraticInputConsistencyFunctionTester}
\leanok
\uses{FABL.NonadaptiveFunctionTester, FABL.affineConsistencyFunctionTester, FABL.padThreeQueryFunctionTester, FABL.quadraticInputConsistencyAddress, FABL.reindexDomainFunctionTester}
Let $n$ be a positive natural number. The associated Boolean tester uses four query
positions and addresses drawn from the disjoint union of the $n$ input positions and
$2^n+2^{n^2}$ proof positions; it is obtained by extending the affine consistency tester
to this address space and padding it to four queries.
\end{definition}

\begin{theorem}[Acceptance probability under proof reindexing]
\lean{FABL.quadraticInputConsistencyFunctionTester_acceptanceProbability}
\label{FABL.quadraticInputConsistencyFunctionTester_acceptanceProbability}
\leanok
\uses{FABL.NonadaptiveFunctionTester.acceptanceProbability, FABL.affineConsistencyFunctionTester, FABL.inputProofWord, FABL.padThreeQueryFunctionTester_acceptanceProbability, FABL.quadraticInputConsistencyAddress, FABL.quadraticInputConsistencyFunctionTester, FABL.quadraticLinearProofTable, FABL.reindexDomainFunctionTester_acceptanceProbability}
\difficulty{4}
Let $n$ be positive, let $w$ be a Boolean input word of length $n$, and let $p$ be a
Boolean proof of length $2^n+2^{n^2}$. The quadratic input consistency tester's
acceptance probability on the combined word $(w,p)$ equals the affine consistency
tester's acceptance probability on $(w,q)$, where $q$ is the length-$2^n$ table obtained
by reading the first $2^n$ proof entries in their indexed order.
\end{theorem}

\begin{theorem}[Sign encoding of the proof table]
\lean{FABL.f₂CubeProofFunction_quadraticLinearProofTable}
\leanok
\uses{FABL.f₂CubeProofFunction, FABL.pcppBoolSignEquiv_eq_signEncode, FABL.quadraticLinearProofFunction, FABL.quadraticLinearProofTable, FABL.signEncode}
\difficulty{3}
For every natural number $n$ and every Boolean table indexed by $2^n+2^{n^2}$ indices,
the function on the binary cube obtained by restricting that table to its first $2^n$
indices and converting its values to signs is pointwise equal to the sign encoding of
the function obtained by the same restriction followed by conversion to $\mathbb{F}_2$ the field
with two elements.
\end{theorem}

\begin{theorem}[Perfect Acceptance of an Encoded Quadratic Proof]
\lean{FABL.quadraticInputConsistencyFunctionTester_complete}
\proofsource{Analysis-of-Boolean-Functions-by-Ryan-ODonnell}{pdf-ffbe9c7b3306-p198-b004}
\label{FABL.quadraticInputConsistencyFunctionTester_complete}
\leanok
\uses{FABL.F₂Cube, FABL.NonadaptiveFunctionTester.acceptanceProbability, FABL.affineConsistencyFunctionTester_acceptanceProbability, FABL.affineCoordinateCorrectionProbability, FABL.encodeHomogeneousQuadraticProof, FABL.f₂CubeProofFunction_quadraticLinearProofTable, FABL.f₂DotProduct, FABL.f₂SignCharacter, FABL.f₂SignCharacter_standardBasis, FABL.f₂StandardBasis, FABL.inputProofWord, FABL.localCorrection, FABL.localCorrection_f₂SignCharacter, FABL.pcppBoolCubeEquiv, FABL.pcppBoolCubeEquiv_apply, FABL.pcppBoolSignEquiv, FABL.pcppBoolSignEquiv_eq_signEncode, FABL.quadraticInputConsistencyFunctionTester, FABL.quadraticInputConsistencyFunctionTester_acceptanceProbability, FABL.quadraticLinearProofFunction_encodeHomogeneousQuadraticProof, FABL.signEncode, FABL.uniformProbability}
\difficulty{4}
For every positive integer $n$ and every Boolean word $w$ of length $n$, the quadratic
input consistency tester accepts with probability $1$ when given the word formed by
concatenating $w$ with the encoded homogeneous quadratic proof of the binary vector
corresponding to $w$.
\end{theorem}

\begin{definition}[Homogeneous quadratic proof tester]
\lean{FABL.homogeneousQuadraticProofFunctionTester}
\statementsource{Analysis-of-Boolean-Functions-by-Ryan-ODonnell}{pdf-ffbe9c7b3306-p198-b003}
\label{FABL.homogeneousQuadraticProofFunctionTester}
\leanok
\uses{FABL.HomogeneousQuadraticSystem, FABL.NonadaptiveFunctionTester, FABL.quadraticInputConsistencyFunctionTester, FABL.quadraticPairProofFunctionTester, FABL.quadraticSystemProofFunctionTester, FABL.uniformMixtureFunctionTester}
Let $n$ be a positive natural number and let $L$ be a finite list of equations in $n$
variables over the field with two elements, each specified by a coefficient matrix and a
right-hand side. The associated Boolean tester uses four query positions and addresses
from the disjoint union of the $n$ input positions and $2^n+2^{n^2}$ proof positions; it
is the uniform mixture of the quadratic-pair proof tester, the quadratic-system proof
tester for $L$, and the quadratic input-consistency tester.
\end{definition}

\begin{definition}[Homogeneous quadratic PCPP system]
\lean{FABL.homogeneousQuadraticPCPPSystem}
\statementsource{Analysis-of-Boolean-Functions-by-Ryan-ODonnell}{pdf-ffbe9c7b3306-p198-b003}
\label{FABL.homogeneousQuadraticPCPPSystem}
\leanok
\uses{FABL.HomogeneousQuadraticSystem, FABL.PCPPSystem, FABL.homogeneousQuadraticProofFunctionTester}
Let $n$ be a positive natural number and let $L$ be a finite list of equations in $n$
variables over the field with two elements, each specified by a coefficient matrix and a
right-hand side. The associated Boolean proof-checking system has proof length
$2^n+2^{n^2}$ and uses four queries; its verifier is the homogeneous quadratic proof
tester for $L$.
\end{definition}

\begin{theorem}[Acceptance probability of a quadratic PCPP system]
\lean{FABL.homogeneousQuadraticPCPPSystem_acceptanceProbability}
\label{FABL.homogeneousQuadraticPCPPSystem_acceptanceProbability}
\leanok
\uses{FABL.HomogeneousQuadraticSystem, FABL.NonadaptiveFunctionTester.acceptanceProbability, FABL.PCPPSystem.acceptanceProbability, FABL.homogeneousQuadraticPCPPSystem, FABL.homogeneousQuadraticProofFunctionTester, FABL.inputProofWord, FABL.quadraticInputConsistencyFunctionTester, FABL.quadraticPairProofFunctionTester, FABL.quadraticSystemProofFunctionTester, FABL.uniformMixtureFunctionTester_acceptanceProbability}
\difficulty{3}
Let $n$ be a positive integer, let $L$ be a finite list of homogeneous quadratic
equations over $\mathbb F_2$ in $n$ variables, let $w\in\{0,1\}^n$ be an input, and let
$p$ be a Boolean proof indexed by the disjoint proof table of size $2^n+2^{n^2}$. The
acceptance probability of the verifier that chooses uniformly among the three component
tests—the quadratic-pair test, the system test, and the input-consistency test—is the
arithmetic mean of the three component tests’ acceptance probabilities on $(w,p)$.
\end{theorem}

\begin{theorem}[Completeness of the quadratic pair proof tester]
\lean{FABL.quadraticPairProofFunctionTester_complete}
\proofsource{Analysis-of-Boolean-Functions-by-Ryan-ODonnell}{pdf-ffbe9c7b3306-p198-b003}
\label{FABL.quadraticPairProofFunctionTester_complete}
\leanok
\uses{FABL.NonadaptiveFunctionTester.acceptanceProbability, FABL.encodeHomogeneousQuadraticProof, FABL.inputProofWord, FABL.pcppBoolCubeEquiv, FABL.quadraticCharacterPairTester_complete, FABL.quadraticLinearProofFunction_encodeHomogeneousQuadraticProof, FABL.quadraticMatrixProofFunction_encodeHomogeneousQuadraticProof, FABL.quadraticPairProofFunctionTester, FABL.quadraticPairProofFunctionTester_acceptanceProbability}
\difficulty{3}
For every Boolean word $w\colon \{1,\ldots,n\}\to\{0,1\}$, the quadratic pair proof
function tester accepts with probability $1$ when given the input word $w$ together with
its encoded homogeneous quadratic proof.
\end{theorem}

\begin{theorem}[Quadratic system proof function tester complete]
\lean{FABL.quadraticSystemProofFunctionTester_complete}
\statementsource{Analysis-of-Boolean-Functions-by-Ryan-ODonnell}{pdf-ffbe9c7b3306-p198-b003}
\label{FABL.quadraticSystemProofFunctionTester_complete}
\leanok
\uses{FABL.HomogeneousQuadraticSystem, FABL.IsBoolHomogeneousQuadraticSystem, FABL.NonadaptiveFunctionTester.acceptanceProbability, FABL.binaryMatrixCharacter_unflattenBinaryMatrix, FABL.encodeHomogeneousQuadraticProof, FABL.flattenBinaryMatrix, FABL.inputProofWord, FABL.linearizedQuadraticSystemCharacterTester_complete, FABL.pcppBoolCubeEquiv, FABL.quadraticMatrixProofFunction_encodeHomogeneousQuadraticProof, FABL.quadraticSystemProofFunctionTester, FABL.quadraticSystemProofFunctionTester_acceptanceProbability, FABL.rankOneSelfMatrix, FABL.satisfiesLinearized_rankOne_iff}
\difficulty{3}
% dataset-pipeline: stub -- prose pending build_informal_statements.py
Statement pending regeneration.
\end{theorem}

\begin{theorem}[Completeness of the quadratic PCPP system]
\lean{FABL.homogeneousQuadraticPCPPSystem_complete}
\proofsource{Analysis-of-Boolean-Functions-by-Ryan-ODonnell}{pdf-ffbe9c7b3306-p198-b003}
\label{FABL.homogeneousQuadraticPCPPSystem_complete}
\leanok
\uses{FABL.HomogeneousQuadraticSystem, FABL.IsBoolHomogeneousQuadraticSystem, FABL.PCPPSystem.acceptanceProbability, FABL.encodeHomogeneousQuadraticProof, FABL.homogeneousQuadraticPCPPSystem, FABL.homogeneousQuadraticPCPPSystem_acceptanceProbability, FABL.pcppBoolCubeEquiv, FABL.quadraticInputConsistencyFunctionTester_complete, FABL.quadraticPairProofFunctionTester_complete, FABL.quadraticSystemProofFunctionTester_complete}
\difficulty{3}
Let $n$ be positive, let $L$ be a finite list of homogeneous quadratic equations over
$\mathbb F_2$ in $n$ variables, and let $w\in\{0,1\}^n$ satisfy every equation in $L$.
There is an encoded proof of length $2^n+2^{n^2}$, formed from the linear and rank-one
quadratic characters determined by the binary encoding of $w$, for which the associated
quadratic PCPP verifier accepts with probability exactly $1$.
\end{theorem}

\begin{definition}[Homogeneous quadratic PCPP rejection rate]
\lean{FABL.homogeneousQuadraticPCPPRejectionRate}
\statementsource{Analysis-of-Boolean-Functions-by-Ryan-ODonnell}{pdf-ffbe9c7b3306-p198-b003}
\label{FABL.homogeneousQuadraticPCPPRejectionRate}
\leanok
For a natural number $n$, define the homogeneous quadratic proof-checking system's
rejection-rate parameter to be $1/100000$.
\end{definition}

\begin{theorem}[Component acceptance from mixture acceptance]
\lean{FABL.homogeneousQuadraticPCPPSystem_component_acceptanceProbability_ge}
\label{FABL.homogeneousQuadraticPCPPSystem_component_acceptanceProbability_ge}
\leanok
\uses{FABL.HomogeneousQuadraticSystem, FABL.NonadaptiveFunctionTester, FABL.NonadaptiveFunctionTester.acceptanceProbability, FABL.PCPPSystem.acceptanceProbability, FABL.component_acceptanceProbability_ge_of_uniformMixture_ge, FABL.homogeneousQuadraticPCPPRejectionRate, FABL.homogeneousQuadraticPCPPSystem, FABL.homogeneousQuadraticProofFunctionTester, FABL.inputProofWord, FABL.quadraticInputConsistencyFunctionTester, FABL.quadraticPairProofFunctionTester, FABL.quadraticSystemProofFunctionTester}
\difficulty{4}
Let $n>0$, let $L$ be a finite system of homogeneous quadratic equations over
$\mathbb{F}_2$, let $\varepsilon\in\mathbb{R}$, and let $w$ and $\pi$ be Boolean words
of lengths $n$ and $2^n+2^{n^2}$, respectively. If the verifier for the associated proof
system accepts $(w,\pi)$ with probability at least $1-\varepsilon/100000$, then each of
its three component testers accepts $(w,\pi)$ with probability at least
$1-(3/100000)\varepsilon$.
\end{theorem}

\begin{theorem}[Character witness from high acceptance]
\lean{FABL.homogeneousQuadraticPCPPSystem_character_witness_of_acceptanceProbability_ge}
\proofsource{Analysis-of-Boolean-Functions-by-Ryan-ODonnell}{
  pdf-ffbe9c7b3306-p198-b001,
  pdf-ffbe9c7b3306-p198-b003
}
\label{FABL.homogeneousQuadraticPCPPSystem_character_witness_of_acceptanceProbability_ge}
\leanok
\uses{FABL.F₂Cube, FABL.HomogeneousQuadraticSystem, FABL.IsClose, FABL.NonadaptiveFunctionTester.acceptanceProbability, FABL.PCPPSystem.acceptanceProbability, FABL.SatisfiesHomogeneousQuadraticSystem, FABL.binaryMatrixCharacter, FABL.binaryMatrixCharacter_unflattenBinaryMatrix, FABL.flattenBinaryMatrix, FABL.flattenBinaryMatrix_unflattenBinaryMatrix, FABL.f₂DotProduct, FABL.homogeneousQuadraticPCPPRejectionRate, FABL.homogeneousQuadraticPCPPSystem, FABL.homogeneousQuadraticPCPPSystem_component_acceptanceProbability_ge, FABL.linearizedQuadraticSystemCharacterTester, FABL.linearizedQuadraticSystemCharacterTester_close_of_acceptanceProbability_ge, FABL.quadraticCharacterOracle, FABL.quadraticCharacterPairRejectionRate, FABL.quadraticCharacterPairTester, FABL.quadraticCharacterPairTester_close_of_acceptanceProbability_ge, FABL.quadraticLinearProofFunction, FABL.quadraticMatrixProofFunction, FABL.quadraticPairProofFunctionTester_acceptanceProbability, FABL.quadraticSystemProofFunctionTester_acceptanceProbability, FABL.rankOneSelfMatrix, FABL.relativeHammingDist, FABL.relativeHammingDist_binaryMatrixCharacter_eq, FABL.relativeHammingDist_triangle_left, FABL.relativeHammingDist_unflattenBinaryMatrix, FABL.satisfiesLinearized_rankOne_iff, FABL.unflattenBinaryMatrix, FABL.𝔽₂}
\difficulty{5}
Let $n$ be positive, let $L$ be a finite list of homogeneous quadratic equations over
$\mathbb F_2$ in $n$ variables, and let $\varepsilon\in[0,1]$. For every Boolean input
$w\in\{0,1\}^n$ and Boolean proof $p$ of length $2^n+2^{n^2}$, if the associated
quadratic PCPP verifier accepts with probability at least $1-\varepsilon/100000$, then
there is a vector $\gamma\in\mathbb F_2^n$ satisfying every equation in $L$ such that
the proof table’s linear part differs from the linear function
$x\mapsto\langle\gamma,x\rangle$ on at most a $(3\varepsilon/1000)$ fraction of inputs.
Its matrix part likewise differs on at most a $(3\varepsilon/1000)$ fraction of matrices
from the function $A\mapsto\langle A,\gamma\gamma^{\mathsf T}\rangle$, where both
pairings are over $\mathbb F_2$.
\end{theorem}

\begin{theorem}[Local correction probability for linear functions]
\lean{FABL.quadraticLinearProof_localCorrectionProbability_ge}
\proofsource{Analysis-of-Boolean-Functions-by-Ryan-ODonnell}{pdf-ffbe9c7b3306-p198-b004}
\label{FABL.quadraticLinearProof_localCorrectionProbability_ge}
\leanok
\uses{FABL.F₂Cube, FABL.IsClose, FABL.affineCoordinateCorrectionProbability, FABL.affineCoordinateCorrectionProbability_ge, FABL.f₂CubeProofFunction, FABL.f₂CubeProofFunction_quadraticLinearProofTable, FABL.f₂DotProduct, FABL.pcppBoolCubeEquiv, FABL.quadraticLinearProofFunction, FABL.quadraticLinearProofTable, FABL.relativeHammingDist, FABL.relativeHammingDist_signEncode, FABL.signEncode}
\difficulty{3}
Let $n$ be a positive integer, let $P$ be a Boolean-valued function on a set of
$2^n+2^{n^2}$ indices, and let $	au:\mathbb{F}_2^n\to\{\pm1\}$ be the sign-valued
function obtained from its first $2^n$ entries. Suppose the Boolean-valued function
encoded by those entries differs from the linear function
$x\mapsto\langle\gamma,x\rangle$ on at most a fraction $\delta$ of inputs, where
$\gamma\in\mathbb{F}_2^n$. Then, for every coordinate $i$, the probability over uniform
$y\in\mathbb{F}_2^n$ that $\tau(y)\tau(y+e_i)$ equals the sign encoding of the $i$th
coordinate of $\gamma$ is at least $1-2\delta$.
\end{theorem}

\begin{theorem}[Closeness to a quadratic solution]
\lean{FABL.homogeneousQuadraticPCPPSystem_close_of_acceptanceProbability_ge}
\proofsource{Analysis-of-Boolean-Functions-by-Ryan-ODonnell}{pdf-ffbe9c7b3306-p198-b003}
\label{FABL.homogeneousQuadraticPCPPSystem_close_of_acceptanceProbability_ge}
\leanok
\uses{FABL.HomogeneousQuadraticSystem, FABL.IsBoolHomogeneousQuadraticSystem, FABL.IsClose, FABL.IsCloseToProperty, FABL.NonadaptiveFunctionTester.acceptanceProbability, FABL.PCPPSystem.acceptanceProbability, FABL.SatisfiesHomogeneousQuadraticSystem, FABL.affineConsistencyFunctionTester, FABL.affineCoordinateCorrectionProbability, FABL.f₂CubeProofFunction, FABL.half_relativeHammingDist_le_affineConsistency_rejection, FABL.homogeneousQuadraticPCPPRejectionRate, FABL.homogeneousQuadraticPCPPSystem, FABL.homogeneousQuadraticPCPPSystem_character_witness_of_acceptanceProbability_ge, FABL.homogeneousQuadraticPCPPSystem_component_acceptanceProbability_ge, FABL.inputProofWord, FABL.pcppBoolCubeEquiv, FABL.quadraticInputConsistencyFunctionTester_acceptanceProbability, FABL.quadraticLinearProofTable, FABL.quadraticLinearProof_localCorrectionProbability_ge}
\difficulty{5}
Let $n$ be positive, let $L$ be a finite list of homogeneous quadratic equations over
$\mathbb F_2$ in $n$ variables, and let $\varepsilon\in[0,1]$. If a Boolean input
$w\in\{0,1\}^n$ and a Boolean proof of length $2^n+2^{n^2}$ make the associated
quadratic PCPP verifier accept with probability at least $1-\varepsilon/100000$, then
there exists a Boolean input satisfying every equation in $L$ whose relative Hamming
distance from $w$ is at most $\varepsilon$.
\end{theorem}

\begin{theorem}[Quadratic system proof verification]
\lean{FABL.homogeneousQuadraticPCPPSystem_isPCPP}
\proofsource{Analysis-of-Boolean-Functions-by-Ryan-ODonnell}{pdf-ffbe9c7b3306-p198-b003}
\label{FABL.homogeneousQuadraticPCPPSystem_isPCPP}
\leanok
\uses{FABL.HomogeneousQuadraticSystem, FABL.IsBoolHomogeneousQuadraticSystem, FABL.IsPCPP, FABL.encodeHomogeneousQuadraticProof, FABL.homogeneousQuadraticPCPPRejectionRate, FABL.homogeneousQuadraticPCPPSystem, FABL.homogeneousQuadraticPCPPSystem_close_of_acceptanceProbability_ge, FABL.homogeneousQuadraticPCPPSystem_complete, FABL.isPCPP_of_exists_close_of_acceptanceProbability_ge, FABL.pcppBoolCubeEquiv}
\difficulty{4}
Let $n>0$ and let $L$ be a finite system of homogeneous quadratic equations in $n$
variables over $\mathbb{F}_2$. Regard its solutions as Boolean words via the
coordinatewise identification of Boolean values with $\mathbb{F}_2$. The associated
proof system has positive rejection rate, is complete (every solution has a proof
accepted with probability $1$), and satisfies the following soundness property: for
every $\varepsilon\in[0,1]$, every Boolean word whose relative Hamming distance from
every solution is greater than $\varepsilon$, and every Boolean proof, the rejection
probability is strictly greater than $\varepsilon/100000$.
\end{theorem}

\begin{definition}[Homogeneous quadratic PCPP description]
\lean{FABL.homogeneousQuadraticPCPPDescription}
\statementsource{Analysis-of-Boolean-Functions-by-Ryan-ODonnell}{
  pdf-ffbe9c7b3306-p198-b003,
  pdf-ffbe9c7b3306-p198-b004
}
\label{FABL.homogeneousQuadraticPCPPDescription}
\leanok
\uses{FABL.HomogeneousQuadraticSystem, FABL.IsBoolHomogeneousQuadraticSystem, FABL.PCPPDescription, FABL.homogeneousQuadraticPCPPRejectionRate, FABL.homogeneousQuadraticPCPPSystem, FABL.homogeneousQuadraticPCPPSystem_isPCPP}
Let $n$ be a positive natural number and let $L$ be a finite list of equations in $n$
variables over the field with two elements, each specified by a coefficient matrix and a
right-hand side. There is a Boolean proof-checking description for the set of Boolean
assignments whose corresponding vectors over the field with two elements satisfy every
equation in $L$: it has proof length $2^n+2^{n^2}$, query count four, and rejection-rate
parameter $1/100000$, and its system satisfies the perfect-completeness and
far-from-satisfying rejection conditions at that rate.
\end{definition}

\begin{theorem}[Zero-dimensional quadratic systems]
\lean{FABL.isBoolHomogeneousQuadraticSystem_zero_universal_or_empty}
\label{FABL.isBoolHomogeneousQuadraticSystem_zero_universal_or_empty}
\leanok
\uses{FABL.HomogeneousQuadraticSystem, FABL.IsBoolHomogeneousQuadraticSystem, FABL.homogeneousQuadraticSystem_zero_universal_or_empty, FABL.pcppBoolCubeEquiv}
\difficulty{4}
For every finite list $L$ of homogeneous quadratic equations over the binary field in
zero variables, either every Boolean assignment on the empty coordinate set satisfies
every equation in $L$, or no such assignment satisfies the system.
\end{theorem}

\begin{theorem}[Homogeneous quadratic pcpp description or zero]
\lean{FABL.homogeneousQuadraticPCPPDescription_or_zero}
\statementsource{Analysis-of-Boolean-Functions-by-Ryan-ODonnell}{pdf-ffbe9c7b3306-p198-b003}
\label{FABL.homogeneousQuadraticPCPPDescription_or_zero}
\leanok
\uses{FABL.HomogeneousQuadraticSystem, FABL.IsBoolHomogeneousQuadraticSystem, FABL.PCPPDescription, FABL.homogeneousQuadraticPCPPDescription, FABL.isBoolHomogeneousQuadraticSystem_zero_universal_or_empty}
\difficulty{4}
% dataset-pipeline: stub -- prose pending build_informal_statements.py
Statement pending regeneration.
\end{theorem}
